\section{Comparación de dos cortes e identificación del límite}\label{app:closing}
Las estimaciones anteriores valen para cada corte admisible fijo. Sus umbrales de \(n\)
pueden depender del corte, por lo que aún hay que justificar el paso al extremo.
La cantidad que aproximan en el tablero inicial es la misma para todos los cortes.
Sean \(a_n=nR_Q(n)\), \(a_*=F(1)-1/12\), \(H=P_T+J-S\), y \(b>0\) un corte
admisible. Existen cortes admisibles arbitrariamente próximos a cero. Definimos a corte positivo
\(I(b)=\int_b^1 k(x)x^{-2}\,dx\) y \(\Phi(b)=a_*+I(b)+H\).
Las tres identificaciones de frontera dan \(\eta(b)=\limsup_n|B_n(b)-H|\to0\).
La identidad interior vale en \([b,1]\). Al combinarla con \eqref{eq:signed-mean},
la expansión de frontera con error \(O(n/j^2)\) y \(n g_T(M_n)\to a_*\), queda
\[
 \limsup_n|a_n-\Phi(b)|\le C_\delta b+\eta(b)=:\varepsilon(b).
\]
\begin{lema}[Comparación de dos cortes]\label{lem:dos-cortes}
Sea \((a_n)\) una sucesión real y, para cada corte admisible \(b\), sea \(\Phi(b)\in\mathbb R\) con
\(\limsup_n|a_n-\Phi(b)|\le\varepsilon(b)\), donde \(\varepsilon(b)\to0\) al bajar el corte. Entonces
\(\Phi(b)\) tiene límite finito \(\Phi_0\) a lo largo de los cortes admisibles y \(a_n\to\Phi_0\).
\end{lema}
\begin{proof}
Para dos cortes admisibles fijos comparamos \(\Phi(b),\Phi(b')\) con el mismo \(a_n\), tomando \(n\)
suficientemente grande para ambas estimaciones. La desigualdad triangular y el límite superior
dan \(|\Phi(b)-\Phi(b')|\le\varepsilon(b)+\varepsilon(b')\). Se cumple el criterio de Cauchy a lo
largo de los cortes admisibles y \(\Phi(b)\to\Phi_0\in\mathbb R\). Para cada corte,
\(\limsup_n|a_n-\Phi_0|\le\varepsilon(b)+|\Phi(b)-\Phi_0|\). Como el lado izquierdo es independiente
de \(b\), al bajar el corte resulta \(a_n\to\Phi_0\).
\end{proof}
En particular \(I(b)\) tiene límite finito a lo largo de dichos cortes. En todo este argumento
el perfil físico, las leyes detenidas y sus estimaciones se han usado a corte positivo fijo.

\begin{proof}[Final de la demostración del Teorema~\ref{teo:principal} y del
 Corolario~\ref{cor:derivada}]
Para identificar el límite con la integral impropia ordinaria queda decidir si \(k(x)/x^2\)
puede tener un término proporcional a \(1/x\). La identidad de orden cero del perfil da
\(k(0)=0\). La resolvente proyectada de líneas es analítica
cerca de cero: el bloque de Schur diagonal tiene la forma \(\Lambda_u+x\mathcal N\), la vacancia
tiene mínimo positivo y el bloque axial tiene brecha espectral. La inversa y las secciones de su
núcleo permanecen acotadas en discos complejos interiores. La integración dominada en las
contracciones finitas de las densidades da así la analiticidad de \(k,G\). Esta regularidad se
obtiene para cualquier valor que pueda tener \(k'(0)\). En un intervalo suficientemente
pequeño \([0,x_0]\), escribimos
\(k(x)=k'(0)x+x^2 h(x)\), con \(h\) analítica y acotada. Para \(0<b<x_0\),
\[
 I(b)=I(x_0)+k'(0)\log(x_0/b)+\int_b^{x_0}h(x)\,dx .
\]
La última integral permanece acotada. La convergencia de \(I(b)\) por los cortes admisibles
fuerza \(k'(0)=0\). De \(k'=d'-xG''\) se deduce \(d'(0)=0\) para el perfil físico.
La singularidad del integrando es entonces evitable y la integral impropia converge
absolutamente. Por consiguiente, el valor de \(\Phi_0\) es
\[
 \lim_{n\to\infty}n\{\log Q(n)-\log Q_{\rm pred}(n)\}
   =F(1)-\tfrac1{12}+\int_0^1\frac{k(x)}{x^2}\,dx+P_T+J-S .
\]
\end{proof}
